\documentclass[a4paper]{article}

\usepackage[top=1in,bottom=1in,left=1in,right=1in]{geometry}
\usepackage[english]{babel}
\usepackage{lastpage} %统计页数
\usepackage{fontspec} %使用非latex字体要使用此包
\usepackage{ulem}
%自定义页眉页脚
\usepackage{fancyhdr} 
\pagestyle{fancy}
\fancyfoot[C]{\fontspec{SimHei} 第{\thepage}页   共\pageref{LastPage}页} 
\renewcommand{\headrulewidth}{0pt}
\renewcommand{\footrulewidth}{0pt}

%更改大列表环境，小列表环境可以在label中改变
\renewcommand{\theenumi}{\Roman{enumi}.}
\renewcommand{\labelenumi}{\theenumi}
\renewcommand{\theenumii}{\Alph{enumii}.}
\renewcommand{\labelenumii}{\theenumii}
\renewcommand{\theenumiii}{\arabic{enumiii}.}
\renewcommand{\labelenumiii}{\theenumiii}
\renewcommand{\theenumiv}{\alph{enumiv}.}
\renewcommand{\labelenumiv}{\theenumiv}
%格式
\setmainfont{Times New Roman} 
\renewcommand{\normalsize}{\fontsize{10.5pt}{\baselineskip}\selectfont}
\renewcommand{\Large}{\fontsize{12pt}{\baselineskip}\selectfont}
\renewcommand{\LARGE}{\fontsize{15.75pt}{\baselineskip}\selectfont}
\setlength{\parindent}{0pt} %首行不缩进
\setlength{\baselineskip}{20pt} %行距

%文档
\begin{document}
	\thispagestyle{fancy}
\includegraphics{header.eps} \\
{\fontspec{SimHei} \Large 一、论文标题} \\
{\fontspec{Arial Black}\Large HJD Type Crane PLC transformation in control circuit \\
	\fontspec{SimHei}\Large HJD型克令吊控制电路的PLC改造} \\
%摘要
{\fontspec{SimHei} \Large 二、摘要与关键词} \\
{\fontspec{Arial Black} \Large Abstract:} \\
\textbf{Background}. HJD type crane is widely used in domestic ships. Traditional HJD type crane adopts relays and contacts for secondary circuit control. The coils and contacts are easily burned after long-term use, frequently causing accidents. \\
\textbf{Aims and approach}. The chief aim of the present work is to provide an available way shifting to PLC control and reduce the use of coils and contacts. PLC, namely, programmable logic controller, is a digital computer used for automation and control. PLC supports multiple analogue and digital inputs and outputs arrangements, with extended temperature ranges, immunity to electrical noise, and resistance to vibration and impact. PLC used in this paper is Siemens S7-200.\\
\textbf{Conclusion}. In this paper, The PLC transformation, including basic protection control process, rose and fall process, brake process, low/medium/high-speed process, has been supplied. There are no failures in experiment after long-term usage. 

{\fontspec{Arial Black} \textbf{\Large Key words: }}HJD, PLC, crane, S7-200, simens.

%大纲
{\fontspec{SimHei} \Large 三、提纲} \\
{\fontspec{Arial Black} \Large Outline} \\
Thesis Sentence: The PLC transformations and the method of implementation, including basic protection control process, rose and fall process, brake process, low/medium/high-speed process, list below. Details can be found in the article.
\begin{enumerate}
\item Basic protection control process is transformed considering overload, overheat, etc.
	\begin{enumerate}
\item Overload protection includes fan motor and crane motor
		\begin{enumerate}
\item Fan motor overload protection is achieved  by fan thermal relay FR2.
\item Crane motor overload protection is achieved  by crane motor thermal relay FR1.
		\end{enumerate}
\item Motor windings overheating protection is achieved by  motor temperature controller ST.
\item Missing phase of power and circuit break protection are achieved by  zero-voltage relay KA1.
\item Emergency forced running is achieved by the contact SB.
	\end{enumerate}
\item Rose and fall process is transformed into rose process and fall process.
	\begin{enumerate}
\item Rose process  is achieved by steering control contact Q0.2.
\item Fall process is  achieved by steering control contact Q0.3.
\item DC delay time relay is required in reversing at high speed.
	\end{enumerate}
\item Brake process is different depending on the conditions of motor.
	\begin{enumerate} 
\item In normal rise or fall state, parking brake   coil pulls in,mechanical parking brake operates.
\item In medium/high gear, DC master switch disconnect, low speed winding connects realizing automatic grade braking.
	\end{enumerate} 
\item Reconstruction of low/medium/high-speed process lists below.
	\begin{enumerate}
\item Reconstruction of low-speed process lists below.
		\begin{enumerate}
\item The crane will run in operational status by releasing rise or fall  contact and the braking contact.
\item The shift from low-speed to medium-speed is achieved by energizing medium-speed winding power.
		\end{enumerate}
\item Reconstruction of medium-speed process lists below.
	\begin{enumerate}
\item Medium speed contact is self-locked and  interlocked with  low/high speed contact.
\item Braking contact and fan contact make sure that the motor will not run at medium speed.
	\end{enumerate}
\item Reconstruction of high-speed process lists below.
	\begin{enumerate}
\item Motor should not run in high speed when heavy loaded.
	\begin{enumerate}
\item The detection of load is achieved  by load relay.
\item The protection of overload is achieved by load contacts.
	\end{enumerate}
\item The diagrams of wiring and control logic must upload to PLC.
	\end{enumerate} 
	\end{enumerate}
	\end{enumerate}

%参考文献
{\fontspec{SimHei} \Large 四、参考文献} \\
{\fontspec{Arial Black} \Large Bibliography} 

\hangafter=1
\setlength{\hangindent}{4em}
Smoczek, J., and J. Szpytko. 2014. ``Evolutionary algorithm-based design of a fuzzy TBF predictive model and TSK fuzzy anti-sway crane control system.''
\textit{Engineering Applications of Artificial Intelligence} 28:190-200. \par

\hangafter=1
\setlength{\hangindent}{4em}
Kłosiński, Jacek. 2005. ``Swing-free stop control of the slewing motion of a mobile crane.''
\textit{Control Engineering Practice} 13:451-460. \par

\hangafter=1
\setlength{\hangindent}{4em}
Kim, Dooroo, and William Singhose. 2010. ``Performance studies of human operators driving double-pendulum bridge cranes.''
\textit{Control Engineering Practice} 18:567-576. \par

\hangafter=1
\setlength{\hangindent}{4em}
Armstrong, N.A., and P.R. 1994. ``MooreA distributed control architecture for intelligent crane automation.''
\textit{Automation in Construction} 3:45-53. \par

\hangafter=1
\setlength{\hangindent}{4em}
Das, S.K., and B.K. Wada. 1990. ``Use of reduced basis technique in the inverse dynamics of large space cranes.''
\textit{Computing Systems in Engineering} 1:577-589. \par

\hangafter=1
\setlength{\hangindent}{4em}
Chang, Chengyuan. 2008. ``Fuzzy projection control law and its application to the overhead crane.''
\textit{Mechatronics} 18:607-615. \par

\hangafter=1
\setlength{\hangindent}{4em}
Armstrong, N.A., and P.R. Moore. ``A distributed control architecture for intelligent crane automation.''
\textit{Automation in Construction} 3:45-53. \par

\hangafter=1
\setlength{\hangindent}{4em}
Rusinski, Eugeniusz, Zaklina Stamboliska, and Przemysław Moczko. 2013. ``Proactive control system of condition of low-speed cement machinery.''
\textit{Automation in Construction} 31:313-324. \par

\hangafter=1
\setlength{\hangindent}{4em}
Gelen, Gökhan, Murat Uzam. 2014. ``The synthesis and PLC implementation of hybrid modular supervisors for real time control of an experimental manufacturing system.''
\textit{Journal of Manufacturing Systems} 535-550. \par

\hangafter=1
\setlength{\hangindent}{4em}
Peng, Xiaohua, Zhilei Cui, Wen Gu, Weiguo Xu, Xuejun Guo. 2014. ``Low level of LAT–PLC-γ1 interaction is associated with Th2 polarized differentiation: A contributing factor to the etiology of asthma.''
\textit{Cellular Immunology} 131-137. \par

\hangafter=1
\setlength{\hangindent}{4em}
Javier Matanza, Sadot Alexandres, Carlos Rodriguez-Morcillo. 2013. ``Performance evaluation of two narrowband PLC systems: PRIME and G3.''
\textit{Computer Standards \& Interfaces} 198-208. \par

%致谢
{\fontspec{SimHei} \Large 五、致谢} \\
{\fontspec{Arial Black} \Large Acknowledgments} \\
The author acknowledges the support of the Shanghai Maritime
University and Merchant Marine college. The author appreciates the valuable
help from instructor Yechun Ling, and express his gratitude to all those who helped him
during the writing of this thesis. The content is solely the responsibility of the author and
does not necessarily represent the official views of the Shanghai Maritime University
and Merchant Marine College.
\end{document}
